%HOMEWORK template for M 325K
\documentclass{article}
\headheight=7pt         \topmargin=14pt
\textheight=574pt       \textwidth=445pt
\oddsidemargin=18pt     \evensidemargin=18pt

%Begin package import

\usepackage{amsmath, amssymb, amsthm}
\usepackage{mathtools}
\usepackage{pdfpages}
\usepackage{tikz-cd}

%End package import
%Begin custom commands

\newcommand{\C}{\mathbb{C}}
\newcommand{\R}{\mathbb{R}}
\newcommand{\Q}{\mathbb{Q}}
\newcommand{\Z}{\mathbb{Z}}
\newcommand{\N}{\mathbb{N}}
\newcommand{\abs}[1]{\lvert #1\rvert}
\newcommand{\RM}[1]{\mathrm{#1}}
\newcommand{\BF}[1]{\textbf{#1}}
\newcommand{\e}{\epsilon}
\newcommand{\ceil}[1]{\lceil{#1}\rceil}
\newcommand{\floor}[1]{\lfloor{#1}\rfloor}
\newcommand{\rarr}[0]{\rightarrow}
\newcommand{\IM}[1]{\text{Im}(#1)}
\newcommand{\Hom}[0]{\text{Hom}}
\newcommand{\Pow}[1]{\text{Pow}(#1)}

%End custom commands
%Begin custom theorems

\newtheorem{innercustomgeneric}{\customgenericname}
\providecommand{\customgenericname}{}
\newcommand{\newcustomtheorem}[2]{%
\newenvironment{#1}[1]
{%
\renewcommand\customgenericname{#2}%
\renewcommand\theinnercustomgeneric{##1}%
\innercustomgeneric%
}
{\endinnercustomgeneric}
}

\newcustomtheorem{THRM}{Theorem}
\newcustomtheorem{LEM}{Lemma}
\newcustomtheorem{EX}{Exercise}
\newcustomtheorem{COR}{Corollary}
\newcustomtheorem{PROB}{Problem}
\newcustomtheorem{claim}{Claim}

\newenvironment{solution}
{\renewcommand\qedsymbol{$\blacksquare$}\begin{proof}[Solution]}
{\end{proof}}

\newenvironment{definition}
{\renewcommand\qedsymbol{$\blacksquare$}\begin{proof}[Definition]}
{\end{proof}}

%End custom theorems
\begin{document}

\title{Normalizers and Sylow Theorems}
\author{Arham Lodha}%put your name here
\date{November 28, 2023}%enter the due date here
\maketitle

\begin{PROB}{1}\label{Problem 1.}
    For finite group G, construct an injective homo $G \to GL_{|G|}(F)$, for any field F. this shows in particular that every finite group occurs as a subgroup of $GL_n(F_p)$ for some p and n. Write down $|GL_n(F_p)|$ and show that the largest power of p that divides this is $k:= 1 + 2 + .. + n-1$.   Look at that number. Exhibit a subgroup of $GL_n(F_p)$ of order $p^k$ (look for a set if  matrices where you have to choose an element of $F_p k$ times).
\end{PROB}
\begin{proof}
    $S_{|G|} \subset GL_{|G|}(F)$ as the permutation matrices. $G \subset S_{|G|}$ because G acting on G by left multiplication. Thus $G \subset GL_{|G|}(F)$. Where $G \to GL_{|G|}(F)$ sends $g$ to its corresponding permutation matrix.

    $|GL_n(F_p)| = \prod_{i = 0}^{i = n - 1}(p^n - p^{i}) = \prod_{i = 0}^{i = n - 1}p^i(p^{n - i} - 1) = \prod_{i = 0}^{i = n - 1}(p^i) \prod_{i = 0}^{i = n - 1}(p^{n - i} - 1) = p^{k}\prod_{i = 0}^{i = n - 1}(p^{n - i} - 1)$ $k := 1 + \ldots + n-1$.

    Matrices which are upper triangular with ones on the diagonal and anything else in the triangle above the diagonal. 
\end{proof}

\bigskip

\begin{PROB}{2}\label{Problem 2.}
    Let H  be a finite group acting on a set X, with $|X|$ not divisible by p. Let $p^k$ be the highest power of p that divides $|H|$. Show that $p^k$ divides the order of some stabilizer. 
\end{PROB}
\begin{proof}
    $|H| = p^kq$ where $p \nmid q$
    .
    X may orbits which divide p and orbits which don't.
    $|X| = kp + |Y|$ where $k$ is the size of the orbits which are divisible by p divided by p and $Y \subseteq X$ is the subset of elements with orbits which don't divide p. $\exists y \in Y$, because $ p \nmid |X|$. 

    $\forall y \in |Y|$, $p \nmid |H \cdot y|$. If $|H \cdot y| = 1$, then $p^kq = |H^y||H\cdot y| = |H^y|$ and so $p^k \mid |H^y|$. If $|H \cdot y| \neq 1$, $|H \cdot y| \mid p^kq$ and since $p$ is prime, $|H \cdot y| \mid q$, let $a := q/|H \cdot y|$. $p^kq = |H^y||H\cdot y| \implies p^ka|H \cdot y| = |H^k||H \cdot y| \implies p^ka = |H^k| \implies p^k \mid |H^k|$
\end{proof}

\bigskip

\begin{PROB}{3a}\label{Problem 3a.}
    Now suppose H is a subgroup of G and G has a p-sylow, P.  Let H act on G/P. Explain why the stabilizer of gP is $ H \cap gPg^{-1}$
\end{PROB}
\begin{proof}
    Suppose G acts on G/P. $\forall gP \in G/P$, $\forall h \in G^{gP}$, $hgP = gP$. Thus $hg = gp \ni p \in P$ thus $h = gpg^{-1}$. Thus $h \in gPg^{-1}$. $\forall h \in gPg^{-1}$, $hgP = gpg^{-1}gP$ for $p \in P$, $gpP = gP = hgP$. Thus $h \in G^{gP}$. Thus $G^{gP} = gPg^{-1}$

    If $H$ acts on $gP$, then stabilizer of $gP$ would be the intersection of $H$ and $G^{gP}$ because $H \leq g$ thus $H^{gP} = H \cap gPg^{-1}$.
\end{proof}
\begin{PROB}{3b}\label{Problem 3b.}
    Show that one of these is a p-sylow subgroup of H.
\end{PROB}
\begin{proof}
    Since $P$ is a p-sylow subgroup of G, $p \nmid [G : P]$. $\exists gP \in G/P \ni p \nmid |H \cdot gP|$ because otherwise $p \mid [G : P]$. By problem 2, we know that $p^k | H^{gP}$ where k is the highest power of p which divides $|H|$. Thus $p \nmid [H : H^{gP}]$
\end{proof}

\bigskip
\begin{PROB}{4a}\label{Problem 4a.}
    Now show that if G is a finite group and P is a sylow p, and H is a p-subgroup (meaning |H| is a power of p), that H is contained in some conjugate of P.
\end{PROB}
\begin{proof}
    By problem 3, $H \cap gPg^{-1}$ is p sylow in H but since $H$ is a p group so $|H| = p^n$ so by definition $|H \cap gPg^{-1}| = p^n$ so $H = H \cap gPg^{-1}$ so $H \subseteq gPg^{-1}$
\end{proof}

\begin{PROB}{4b}\label{Problem 4b.}
    In particular conclude that all the sylow p subgroups are conjugate. 
\end{PROB}
\begin{proof}
    Suppose $H$ is p sylow it has order $p^k$ and all P sylow subgroups of G have order $p^k$. Since $H \subseteq gPg^{-1}$ and since $|H| = p^k$ and $|gPg^{-1}| = p^k$. Thus $H = gPg^{-1}$. 
\end{proof}

\bigskip

\begin{PROB}{6.4}\label{Problem 6.4.}
    Let H be a normal subgroup of prime order p in a finite group G. Suppose that p is the smallest prime that divides the order of G. Prove that H is in the center Z(G).
\end{PROB}
\begin{proof}
    Let G act on $H \setminus {e}$ by conjugation. That gives a homomorphism of $f: G \to S_{p-1}$. $f_*(G)$ is a subgroup of $S_{p-1}$ and thus divides $S_{p-1}$ but $f_*(G) | G$ and p is the smallest prime which divides the order of G thus $f_*{G} = 1, p$. If $|f_*{G}| = p$ then it can't divide $S_{p-1}$ so $|f_*{G}| = 1$. Thus everything in G when it conjugates H, $ghg^{-1} = h$ thus everything in g commutes with H. Thus $H \in Z(G)$.
\end{proof}

\bigskip

\begin{PROB}{6.6}\label{Problem 6.6.}
    Let H be a proper subgroup of a finite group G. Prove: The group G is not the union of the conjugate subgroups of H. There is a conjugacy class C that is disjoint from H.
\end{PROB}
\begin{proof}
    The number of conjugate subgroups of H in G is [G:N(H)]=|G|/|N(H)|. But $|N(H)|\geq|H|$, so the number of conjugate subgroups is at most |G|/|H|. However, all conjugate subgroups share the identity element, so the size of the union of these subgroups is at most $1+(|H| - 1)|G|/|H|=|G|+1 - |G|/|H|$. Since |H| is a proper subgroup, this value is less than |G| , so this proves part (a). Part (b) is obviously equivalent to part (a).
\end{proof}

\begin{PROB}{7.3}\label{Problem 7.3.}
    How many elements of order 5 might be contained in a group of order 20?
\end{PROB}
\begin{proof}
    $|G| = 20 = 5 * 4$. All p sylow groups are conjugate. G acts on the sylow p groups by conjugation and the set of sylow p groups is transitive. So number of sylow p groups must divide 20. By 3rd sylow theorems, n = number of sylow p groups $n \equiv 1 \mod p$.Thus n must divide 20 and must be congruent to 1 mod 5. Thus there is only 1 sylow p group and thus only 1 subgroup with order 5. Every element except e will have order 5. Thus 4 elements have order 5. 
\end{proof}

\bigskip

\begin{PROB}{7.4a}\label{Problem 7.4a.}
    Prove that no simple group has order pq, where p and q are prime.
\end{PROB}
\begin{proof}
    Assume the contrary, let G be a simple group with order pq where p and q are prime where q > p. Then there exists a q sylow group P with order q. The n which number of q sylow subgroups must divide pq and be congruent to 1 mod q, then n must divide p thus $n = 1$ or p. But since $q > p$, if $n = p$ $p \equiv 1 mod q$ which would mean $p = 1$ but p is prime so $n = 1$. Thus the number of conjugate subgroups is 1 and thus q sylow group is normal. Thus a contradiction forms and G is no longer normal. Thus no simple group have order pq.  
\end{proof}
\bigskip
\begin{PROB}{7.4b}\label{Problem 7.4b.}
    Prove that no simple group has order $p^2q$, where p and q are prime.
\end{PROB}
\begin{proof}
    There are $1, p, p^2$ q sylow subgroups.

    \begin{enumerate}{}
        \item If there are only 1 q-sylow subgroup, then it is normal. 
        \item If there are p q-sylow subgroups. Then $p \equiv 1 \mod q \implies p > q$. Let n be the number of $p^2$-sylow subgroups. $n \mod p = 1$ and $n | p^2q$, so $n = 1, q$ since q is prime. If $n = q$ then $q \mod p = 1$ which is impossible since $p > q$ and q is prime. Thus $n =1$. Thus the $p^2$ - sylow subgroup is normal.
        \item All the q sylow subgroups are cyclic and thus don't share any element except the identity. Thus the number of elements with order q is $p^2(q - 1)$. Any element in $p^2$ sylow subgroup doesn't have order q. The number of elements in G without order q is $p^2q - p^2(q - 1) = p^2$. Thus there is only 1 $p^2$ sylow subgroup which is normal.
    \end{enumerate}
\end{proof}

\bigskip

\begin{PROB}{7.5a}\label{Problem 7.5a.}
    Find Sylow 2-subgroups for $D_{10}$
\end{PROB}
\begin{proof}
    $D_{10} = \langle x, y | x^{10} = y^2 = e, yxy^{-1} = x^{-1} \rangle$
    One sylow 2-subgroup is $\{e, x^5, y, x^5y\}$. All conjugations are of the form $\{e, x^5, x^ky, x^{k + 5}y\}$
\end{proof}

\begin{PROB}{7.5b}\label{Problem 7.5b.}
    Find Sylow 2-subgroups for $T$
\end{PROB}
\begin{proof}
    Identity and the 3 180 degree rotations. This is Klien 4 subgroup which is normal and thus is the only 2-sylow
\end{proof}

\begin{PROB}{7.5c}\label{Problem 7.5c.}
    Find Sylow 2-subgroups for $O$
\end{PROB}
\begin{proof}
    The group O has order 24 = 8 * 3 = $2^3 * 3$, so we seek a subgroup of order 8. $O \cong C$. $D_4 < O \cong S_4$. There are 3 $D_4$ subgroups in $S_4$ because $D_4$ is determined by the 4 elements in the set with a adjacency relation and there are 3 possible adjacency relation. Because for any point there is 3 elements in the set which it is not adjacent to and this fixes the rest. Thus there are 3 $2$ sylow subgroups each which are $D_4$.
\end{proof}

\bigskip


\begin{PROB}{7.5c}\label{Problem 7.5c.}
    Find Sylow 2-subgroups for $I$
\end{PROB}
\begin{proof}
    The group I has order 60, so we seek a subgroup of order 4. The Klein 4 subgroup in $I \cong A_5$ is a sylow subgroup. Furthermore there is only 1 as this is only 4 element subgroup of $A_5$.
\end{proof}

\bigskip

\begin{PROB}{7.5.6}\label{Problem 7.5.6.}
    
\end{PROB}
\begin{proof}
    According to Proposition 7.7.7(c), such a group has an element x of order 7 and an element y of order 3, and these elements satisfy $yxy^{-1}=x$ . Now the only elements of $S_7$ that have order 7 are the 7-cycles, so, without loss of generality, let x=(1 2 3 4 5 6 7) , so that x2=(1 3 5 7 2 4 6) . We must find y such that y relabels x to $x^2$ and has order 3. The obvious relabelling from the way x and $x^2$ are written, above, is y=(2 3 5)(4 7 6) , which indeed has order 3. Since we have the rule yx=$x^2$y , the elements $x^iy^j$ with i from 0 to 6 and j from 0 to 2 form a subgroup of S7 of order at most 21. This subgroup contains both ⟨x⟩ and ⟨y⟩ as subgroups, so its order is divisible by 21. Therefore its order is exactly 21.
\end{proof}


\end{document}